\documentclass[11pt,a4paper]{scrartcl}


\usepackage[utf8]{inputenc}
\usepackage[T1]{fontenc}
\newcommand{\ats}[1]{\par\noindent\textit{Atsakymas:} #1}
\usepackage{cancel}
\usepackage{tikz}
\usepackage{lmodern} 
\usepackage{amsmath,amssymb,amsthm}
\usepackage{graphicx}
\usepackage[dvipsnames,svgnames]{xcolor}
\usepackage{hyperref}
\usepackage{mathrsfs}
\usepackage[shortlabels]{enumitem}
\usepackage[obeyFinal,textsize=scriptsize,shadow,loadshadowlibrary]{todonotes}
\usepackage{mathtools}
\usepackage{microtype}

%%% Paragraph layout
\setlength{\parindent}{0pt}
\setlength{\parskip}{0.6\baselineskip}

%%% Colored sections
\renewcommand*{\sectionformat}%
  {\color{purple}\S\thesection\enskip}
\renewcommand*{\subsectionformat}%
  {\color{purple}\S\thesubsection\enskip}
\renewcommand*{\subsubsectionformat}%
  {\color{purple}\S\thesubsubsection\enskip}
\addtokomafont{paragraph}{\color{orange!35!black}\P\ }
\KOMAoptions{numbers=noenddot}

%%% Title
\addtokomafont{subtitle}{\Large}
\setkomafont{author}{\Large\scshape}
\setkomafont{date}{\Large\normalsize}
% Neigiamas vertikalus tarpas, kuris pakelia dokumento antraštę į viršų. 
% Galite keisti -2.5cm į -3cm (aukščiau) arba -1.5cm (žemiau).
\titlehead{\vspace*{-7.5cm}} 

% PRIDĖTA: Automatiškai perrašome datą į mokyklos pavadinimą
\AtBeginDocument{%
  \date{\normalsize Vilniaus licėjus}
}

%%% Page setup
\usepackage[headsepline]{scrlayer-scrpage}
\renewcommand{\headfont}{}
\addtolength{\textheight}{3.14cm}
\setlength{\footskip}{0.5in}
\setlength{\headsep}{10pt}
% Puslapio antraštė turi rodyti tik konkrečius dokumento duomenis.
% Nustatome juos aiškiai, kad neatsirastų "\@author" / "\@title" arba senų reikšmių.
\makeatletter
\AtBeginDocument{%
  \ihead{\footnotesize\textbf{Andrius Berniukevičius}}%
  \ohead{\footnotesize\textbf{\@title}}%
}
\makeatother
\automark{section}
\chead{}
\cfoot{\pagemark}

%%% Macros
\providecommand{\ol}{\overline}
\providecommand{\ul}{\underline}
\providecommand{\wt}{\widetilde}
\providecommand{\wh}{\widehat}
\providecommand{\eps}{\varepsilon}
\providecommand{\half}{\frac{1}{2}}
\providecommand{\inv}{^{-1}}
\newcommand{\dang}{\measuredangle} %% Directed angle
\newcommand{\vocab}[1]{\textbf{\color{blue}\sffamily #1}}
\providecommand{\CC}{\mathbb C}
\providecommand{\FF}{\mathbb F}
\providecommand{\NN}{\mathbb N}
\providecommand{\QQ}{\mathbb Q}
\providecommand{\RR}{\mathbb R}
\providecommand{\ZZ}{\mathbb Z}
\providecommand{\ts}{\textsuperscript}
\providecommand{\dg}{^\circ}
\providecommand{\ii}{\item}
\providecommand{\defeq}{\coloneqq}
\DeclareMathOperator*{\lcm}{lcm}
\DeclareMathOperator*{\argmin}{arg min}
\DeclareMathOperator*{\argmax}{arg max}
\providecommand{\hrulebar}{\par
\hspace{\fill}\rule{0.95\linewidth}{.7pt}\hspace{\fill}
\par\nointerlineskip \vspace{\baselineskip}}

%%% Optional colored theorems
\usepackage{thmtools}
\usepackage[framemethod=TikZ]{mdframed}
\mdfdefinestyle{mdbluebox}{%
  roundcorner=10pt,
  linewidth=1pt,
  skipabove=12pt,
  innerbottommargin=9pt,
  skipbelow=2pt,
  linecolor=blue,
  nobreak=true,
  backgroundcolor=TealBlue!5,
}
\declaretheoremstyle[
  headfont=\sffamily\bfseries\color{MidnightBlue},
  mdframed={style=mdbluebox},
  headpunct={\\[3pt]},
  postheadspace={0pt}
]{thmbluebox}
\mdfdefinestyle{mdredbox}{%
  linewidth=0.5pt,
  skipabove=12pt,
  frametitleaboveskip=5pt,
  frametitlebelowskip=0pt,
  skipbelow=2pt,
  frametitlefont=\bfseries,
  innertopmargin=4pt,
  innerbottommargin=8pt,
  nobreak=true,
  backgroundcolor=Salmon!5,
  linecolor=RawSienna,
}
\declaretheoremstyle[
  headfont=\bfseries\color{RawSienna},
  mdframed={style=mdredbox},
  headpunct={\\[3pt]},
  postheadspace={0pt},
]{thmredbox}
\mdfdefinestyle{mdgreenbox}{%
  skipabove=8pt,
  linewidth=2pt,
  rightline=false,
  leftline=true,
  topline=false,
  bottomline=false,
  linecolor=ForestGreen,
  backgroundcolor=ForestGreen!5,
}
\declaretheoremstyle[
  headfont=\bfseries\sffamily\color{ForestGreen!70!black},
  bodyfont=\normalfont,
  spaceabove=2pt,
  spacebelow=1pt,
  mdframed={style=mdgreenbox},
  headpunct={ --- },
]{thmgreenbox}
\mdfdefinestyle{mdblackbox}{%
  skipabove=8pt,
  linewidth=3pt,
  rightline=false,
  leftline=true,
  topline=false,
  bottomline=false,
  linecolor=black,
  backgroundcolor=RedViolet!5!gray!5,
}
\declaretheoremstyle[
  headfont=\bfseries,
  bodyfont=\normalfont\small,
  spaceabove=0pt,
  spacebelow=0pt,
  mdframed={style=mdblackbox}
]{thmblackbox}

%% Aplinkos su rėmeliais (thmtools)
\declaretheorem[style=thmbluebox,name=Teorema]{theorem}
\declaretheorem[style=thmbluebox,name=Lema]{lemma}
\declaretheorem[style=thmbluebox,name=Savybė]{property}
\declaretheorem[style=thmbluebox,name=Teiginys]{proposition}
\declaretheorem[style=thmbluebox,name=Išvada]{corollary}
\declaretheorem[style=thmbluebox,name=Teorema,numbered=no]{theorem*}
\declaretheorem[style=thmbluebox,name=Lema,numbered=no]{lemma*}
\declaretheorem[style=thmbluebox,name=Teiginys,numbered=no]{proposition*}
\declaretheorem[style=thmbluebox,name=Išvada,numbered=no]{corollary*}
\declaretheorem[style=thmgreenbox,name=Algoritmas]{algorithm}
\declaretheorem[style=thmgreenbox,name=Algoritmas,numbered=no]{algorithm*}
\declaretheorem[style=thmgreenbox,name=Teiginys]{claim}
\declaretheorem[style=thmgreenbox,name=Teiginys,numbered=no]{claim*}
\declaretheorem[style=thmredbox,name=Pavyzdys]{example}
\declaretheorem[style=thmredbox,name=Pavyzdys,numbered=no]{example*}
\declaretheorem[style=thmblackbox,name=Pastaba]{remark}
\declaretheorem[style=thmblackbox,name=Pastaba,numbered=no]{remark*}

%% Standartinės amsthm aplinkos (Apibrėžimai ir užduotys)
\theoremstyle{definition}
\ifcsname conjecture\endcsname\else\newtheorem{conjecture}{Hipotezė}\fi
\ifcsname definition\endcsname\else\newtheorem{definition}{Apibrėžimas}\fi
\ifcsname fact\endcsname\else\newtheorem{fact}{Faktas}\fi
\ifcsname answer\endcsname\else\newtheorem{answer}{Atsakymas}\fi
\ifcsname ques\endcsname\else\newtheorem{ques}{Klausimas}\fi
\ifcsname exercise\endcsname\else\newtheorem{exercise}{Užduotis}\fi
\ifcsname problem\endcsname\else\newtheorem{problem}{Uždavinys}[section]\fi
\ifcsname conjecture*\endcsname\else\newtheorem*{conjecture*}{Hipotezė}\fi
\ifcsname definition*\endcsname\else\newtheorem*{definition*}{Apibrėžimas}\fi
\ifcsname fact*\endcsname\else\newtheorem*{fact*}{Faktas}\fi
\ifcsname answer*\endcsname\else\newtheorem*{answer*}{Atsakymas}\fi
\ifcsname ques*\endcsname\else\newtheorem*{ques*}{Klausimas}\fi
\ifcsname exercise*\endcsname\else\newtheorem*{exercise*}{Užduotis}\fi
\ifcsname problem*\endcsname\else\newtheorem*{problem*}{Uždavinys}\fi

\newcounter{answerssection}
\newcommand{\answerssection}{%
  \setcounter{answerssection}{\value{section}}%
  \section{Atsakymai}%
  \setcounter{problem}{0}%
  \renewcommand{\theproblem}{\arabic{answerssection}.\arabic{problem}}%
}

%% GALUTINIS SPRENDIMAS PROOF APLINKAI
%% --- PRADŽIA: Įrodymo aplinkos sutvarkymas ---
\makeatletter

% 1. Užtikriname, kad žodis būtų "Įrodymas", net jei "babel" paketas bando jį perrašyti
\AtBeginDocument{%
  \renewcommand{\proofname}{Įrodymas}%
  \@ifpackageloaded{babel}{%
    \addto\captionslithuanian{\renewcommand{\proofname}{Įrodymas}}%
  }{}%
}

% 2. Priverstinai pašaliname scrartcl klasės bold šriftą ir paliekame TIK italic
% 3. Sujungiame su tuščiu kvadratėliu dešiniajame kampe.
\renewcommand{\qedsymbol}{\ensuremath{\Box}}
\renewenvironment{proof}[1][\proofname]{\par
  \pushQED{\qed}%
  \normalfont \topsep6\p@\@plus6\p@\relax
  \trivlist
  \item[\hskip\labelsep
        \normalfont\itshape % Priverčia naudoti tik pasvirą šriftą, kaip nuotraukoje
    #1\@addpunct{.}]\ignorespaces
}{%
  \popQED\endtrivlist\@endpefalse
}
\newenvironment{solution}{\par
  \normalfont \topsep6\p@\@plus6\p@\relax
  \trivlist
  \item[\hskip\labelsep
        \normalfont\itshape
    \textit{Sprendimas}\@addpunct{.}]\ignorespaces
}{%
  \endtrivlist\@endpefalse
}
\newcommand{\ans}[1]{\par\noindent\textit{Atsakymas:} #1}
\makeatother


\title{Algebrinės galios mokymai\\ I dalis}
\author{Anton Anakin Liutvinas ir Pijus Palpatinas Piekus}

\newenvironment{solution}
    {\begin{list}{}{\setlength{\leftmargin}{10pt}\setlength{\rightmargin}{10pt}}\item[]}
    {\end{list}}

\begin{document}

\maketitle

\section{Naudingos algebrinės savybės}

\begin{enumerate}
    \renewcommand{\labelenumi}{\textbf{\theenumi.}}

    \item $(x\pm y)^2=x^2\pm 2xy+y^2$.
    \item $x^2-y^2=(x-y)(x+y)$.
    \item $(x\pm y)^3=x^3\pm 3x^2y+3xy^2\pm y^3$.
    \item $x^3\pm y^3=(x\pm y)(x^2\mp xy+y^2)$.
    \item $x^3+y^3+z^3=(x+y+z)(x^2+y^2+z^2-xy-xz-yz)+3xyz$.
    \item $x^4+x^2y^2+y^4=(x^2+xy+y^2)(x^2-xy+y^2)$.
    \item Jei $x_1,x_2,\ldots,x_n\in\mathbb R$, tai $x_1^2+x_2^2+\cdots+x_n^2=0\Longrightarrow x_1=x_2=\cdots=x_n=0$.
\end{enumerate}


\section{Lygčių sistemos}

\subsection{Teorija}

Pamačius lygčių sistemą verta pirmiausia ieškoti ryšio tarp jos lygčių.
Dažnai padeda:

\begin{itemize}
    \item sudėti arba atimti lygtis;
    \item padauginti lygtis arba vieną lygtį iš tinkamo skaičiaus;
    \item dalyti lygtis ar reiškinius tik tada, kai žinome, kad daliklis nėra lygus nuliui;
    \item ieškoti simetrijos;
    \item sudaryti pilnuosius kvadratus;
    \item panaudoti faktą, kad realiojo skaičiaus kvadratas yra neneigiamas.
\end{itemize}

Svarbu: vien tai, kad kintamųjų yra daugiau negu lygčių, dar nereiškia,
kad sistemą būtina spręsti nelygybėmis arba kad ji turi be galo daug
sprendinių. Reikia nagrinėti konkrečią sistemos struktūrą.

\begin{example}

Išspręskite lygčių sistemą
\[
\left\{
\begin{array}{l}
xy=-30,\\
yz=-12,\\
zx=10.
\end{array}
\right.
\]
\end{example}

\begin{solution}
\textbf{Sprendimas}

Kadangi
\[
xy\neq0,\qquad yz\neq0,\qquad zx\neq0,
\]
tai
\[
x\neq0,\qquad y\neq0,\qquad z\neq0.
\]

Sudauginame visas tris lygtis:
\[
(xy)(yz)(zx)=(-30)(-12)\cdot10.
\]

Kairėje pusėje gauname
\[
x^2y^2z^2=(xyz)^2,
\]
todėl
\[
(xyz)^2=3600=60^2.
\]

Vadinasi,
\[
xyz=60
\quad\text{arba}\quad
xyz=-60.
\]

\textbf{1 atvejis:}
\[
xyz=60.
\]

Tuomet
\[
x=\frac{xyz}{yz}
=\frac{60}{-12}
=-5,
\]

\[
y=\frac{xyz}{zx}
=\frac{60}{10}
=6,
\]

\[
z=\frac{xyz}{xy}
=\frac{60}{-30}
=-2.
\]

Gauname
\[
(x,y,z)=(-5,6,-2).
\]

\textbf{2 atvejis:}
\[
xyz=-60.
\]

Analogiškai,
\[
x=5,\qquad y=-6,\qquad z=2.
\]

Abu trejetai tenkina pradines lygtis.

\textbf{Atsakymas:} $(-5;6;-2)$ ir $(5;-6;2)$.
\end{solution}

\begin{example}

Raskite visus realiųjų skaičių trejetus $(x,y,z)$, tenkinančius sistemą

\[
\left\{
\begin{array}{l}
x^2+9=4y,\\
y^2+1=6z,\\
z^2+4=2x.
\end{array}
\right.
\]
\end{example}

\begin{solution}
\textbf{Sprendimas}

Sudėkime visas tris lygtis:
\[
x^2+y^2+z^2+14=2x+4y+6z.
\]

Visus narius perkeliame į kairę pusę:
\[
x^2-2x+y^2-4y+z^2-6z+14=0.
\]

Išskiriame pilnuosius kvadratus:
\[
(x^2-2x+1)
+
(y^2-4y+4)
+
(z^2-6z+9)
=0.
\]

Taigi
\[
(x-1)^2+(y-2)^2+(z-3)^2=0.
\]

Kadangi nagrinėjame realiuosius skaičius, kiekvienas kvadratas yra
neneigiamas. Todėl jų suma gali būti lygi nuliui tik tada, kai

\[
x-1=0,\qquad y-2=0,\qquad z-3=0.
\]

Vadinasi,
\[
x=1,\qquad y=2,\qquad z=3.
\]

Tačiau tai dar tik galimas sprendinys. Patikriname jį pradinėje sistemoje.

Pirmojoje lygtyje:
\[
1^2+9=10,
\]
bet
\[
4\cdot2=8.
\]

Lygybė negalioja, todėl trejetas $(1,2,3)$ nėra pradinės sistemos
sprendinys.

\textbf{Atsakymas:} sprendinių nėra.
\end{solution}

\subsection{Uždaviniai}

\begin{enumerate}
    \renewcommand{\labelenumi}{\textbf{\theenumi.}}

    \item Raskite visus sprendinius:
    \[
    \left\{
    \begin{array}{l}
    a+b+c+d=4,\\
    a+b+c+e=8,\\
    a+b+d+e=12,\\
    a+c+d+e=16,\\
    b+c+d+e=20.
    \end{array}
    \right.
    \]

    \item Išspręskite lygčių sistemą:
    \[
    \left\{
    \begin{array}{l}
    2x^2-2x+2=y+z,\\
    2y^2-2y+2=z+x,\\
    2z^2-2z+2=x+y.
    \end{array}
    \right.
    \]

    \item Realieji skaičiai $x$ ir $y$ tenkina lygybes
    \[
    xy=10
    \]
    ir
    \[
    (x+1)(y+1)=20.
    \]
    Kam lygi reiškinio
    \[
    (x+2)(y+2)
    \]
    reikšmė?

    \item Su kuriomis realiosiomis $a$ reikšmėmis lygtys
    \[
    x^3+ax+1=0
    \]
    ir
    \[
    x^4+ax^2+1=0
    \]
    turi bendrą šaknį?

    \item Išspręskite lygčių sistemą:
    \[
    \left\{
    \begin{array}{l}
    x+xy+x^2=9,\\
    y+xy+y^2=-3.
    \end{array}
    \right.
    \]

    \item Duoti teigiami realieji skaičiai $x,y,z$, kurie tenkina
    \[
    \left\{
    \begin{array}{l}
    x+\dfrac1y=4,\\[0.2cm]
    y+\dfrac1z=1,\\[0.2cm]
    z+\dfrac1x=\dfrac73.
    \end{array}
    \right.
    \]
    Raskite $xyz$.

    \item Raskite visus sprendinius:
    \[
    \left\{
    \begin{array}{l}
    \dfrac{x}{y}+\dfrac{y}{z}+\dfrac{z}{x}=3,\\[0.3cm]
    \dfrac{y}{x}+\dfrac{z}{y}+\dfrac{x}{z}=3,\\[0.3cm]
    \dfrac1x+\dfrac1y+\dfrac1z=2.
    \end{array}
    \right.
    \]

    \item Raskite visus realiuosius skaičius $x$, $y$ ir $z$, kurie tenkina
    \[
    x^2+y^2+z^2=x-z=2.
    \]

    \item Išspręskite lygčių sistemą
    \[
    \left\{
    \begin{array}{l}
    x^2+4y^2-2z^4=18,\\
    x+2y+2z^2=40,\\
    z^4-2xy=9.
    \end{array}
    \right.
    \]

\end{enumerate}


\section{Keitiniai}

\subsection{Teorija}

Olimpiadiniuose uždaviniuose dažnai pasitaiko lygčių, kurios iš pirmo
žvilgsnio atrodo sudėtingos, tačiau jose kartojasi tas pats reiškinys.
Tokiu atveju naudinga tą reiškinį pakeisti nauju kintamuoju.

Pavyzdžiui, galima įsivesti keitinius

\[
t=x+y,\qquad
t=x-y,\qquad
t=xy,
\]

\[
t=a^x,\qquad
t=\sqrt{f(x)},\qquad
t=\frac1x,\qquad
t=\frac{x}{y}.
\]

Kartais patogu naudoti ir kelis keitinius vienu metu, pavyzdžiui,

\[
s=x+y,\qquad p=xy.
\]

Keitinio tikslas -- sudėtingą lygtį paversti paprastesne, dažnai
kvadratine arba kita lengvai išskaidoma lygtimi.

Atlikus keitinį būtina:

\begin{enumerate}
    \item atsižvelgti į pradinės lygties apibrėžimo sritį;
    \item išspręsti lygtį su naujuoju kintamuoju;
    \item grįžti prie pradinio kintamojo;
    \item patikrinti gautus sprendinius, jei atlikti pertvarkymai galėjo sukurti pašalinių šaknų.
\end{enumerate}

\begin{example}

Išspręskite lygtį
\[
4^x+8=6\cdot2^x.
\]
\end{example}

\begin{solution}
\textbf{Sprendimas}

Pastebime, kad
\[
4^x=(2^2)^x=(2^x)^2.
\]

Todėl įsiveskime keitinį
\[
t=2^x.
\]

Kadangi
\[
2^x>0,
\]
turime
\[
t>0.
\]

Pradinė lygtis tampa
\[
t^2+8=6t.
\]

Perkeliame visus narius į vieną pusę:
\[
t^2-6t+8=0.
\]

Išskaidome dauginamaisiais:
\[
(t-2)(t-4)=0.
\]

Taigi
\[
t=2
\quad\text{arba}\quad
t=4.
\]

Grįžtame prie pradinio kintamojo.

Jei
\[
2^x=2,
\]
tai
\[
x=1.
\]

Jei
\[
2^x=4=2^2,
\]
tai
\[
x=2.
\]

\textbf{Atsakymas:} $x=1$ arba $x=2$.
\end{solution}

\begin{example}

Tegul
\[
S=(x-1)^4+4(x-1)^3+6(x-1)^2+4(x-1)+1.
\]
Kam lygus $S$?
\end{example}

\begin{solution}
\textbf{Sprendimas}

Visuose nariuose kartojasi reiškinys
\[
x-1.
\]

Todėl įsiveskime keitinį
\[
t=x-1.
\]

Tuomet
\[
S=t^4+4t^3+6t^2+4t+1.
\]

Pastebime Niutono binomo formulę:
\[
(t+1)^4=t^4+4t^3+6t^2+4t+1.
\]

Vadinasi,
\[
S=(t+1)^4.
\]

Grįžtame prie $x$:
\[
t=x-1.
\]

Todėl
\[
S=(x-1+1)^4=x^4.
\]

\textbf{Atsakymas:} $x^4$.
\end{solution}

\subsection{Uždaviniai}

\begin{enumerate}
\renewcommand{\labelenumi}{\textbf{\theenumi.}}

    \item Supaprastinkite reiškinį
    \[
    A=(t-1)^3+3(t-1)^2+3(t-1)+1.
    \]

    \item Realieji skaičiai $x$ ir $y$ tenkina lygybę
    \[
    x^2+2xy+y^2+x+y-2=0.
    \]
    Raskite visas reiškinio $x+y$ reikšmes.

    \item Išspręskite lygtį
    \[
    x^2+18x+30
    =
    2\sqrt{x^2+18x+45}.
    \]

    \item Išspręskite lygtį realiųjų skaičių aibėje:
    \[
    100^{1/x}+25^{1/x}
    =
    4{,}25\cdot50^{1/x},
    \qquad x\ne0.
    \]

    \item Išspręskite lygtį realiųjų skaičių aibėje:
    \[
    y^2+5y-\frac{36}{y^2+5y}=0.
    \]

    \item Išspręskite lygtį
    \[
    2^x+3^x-4^x+6^x-9^x=1.
    \]

    \item Išspręskite lygtį
    \[
    \frac{1}{x(x+8)}
    -
    \frac{1}{(x+4)^2}
    =
    \frac49.
    \]

    \item Išspręskite lygtį
    \[
    \sqrt{x^2-1}
    =
    x\sqrt{x}-\sqrt{x-1}.
    \]

    \item Išspręskite lygtį
    \[
    \frac{x^2+11x-6}{x^2-x-6}
    =
    \frac{26x}{x^2-6}.
    \]

    \item Išspręskite lygtį
    \[
    \frac{8^x+27^x}{12^x+18^x}
    =
    \frac76.
    \]

    \item Išspręskite lygtį
    \[
    \frac{(3-\sqrt{x})^4}{\sqrt{x}}
    +
    \frac{x^2}{3-\sqrt{x}}
    =
    \frac{33}{2}.
    \]

    \item Išspręskite lygtį
    \[
    2x^2+6x+9
    =
    7x\sqrt{2x+3}.
    \]

    \item Išspręskite lygtį
    \[
    \frac{5x-4}{x^2}
    +
    \frac{70}{(x+2)^2}
    =
    1.
    \]

    \item Išspręskite lygtį realiaisiais skaičiais:
    \[
    \sqrt[4]{x+32}
    -
    \sqrt[4]{x-33}
    =
    1.
    \]

    \item Išspręskite lygtį
    \[
    x
    =
    11+
    \sqrt{
    \frac{5x-56}{(x-13)(x-14)}
    }.
    \]

\end{enumerate}


\section{Atsakymai}

\subsection{Lygčių sistemos}

\begin{enumerate}
    \renewcommand{\labelenumi}{\textbf{\theenumi.}}

    \item $(a,b,c,d,e)=(-5,-1,3,7,11)$.
    \item $(x,y,z)=(1,1,1)$.
    \item $32$.
    \item $a=-2$.
    \item $(x,y)=(3,-1)$ arba $(x,y)=\left(-\frac92,\frac32\right)$.
    \item $xyz=1$.
    \item $(x,y,z)=\left(\frac32,\frac32,\frac32\right)$.
    \item $(x,y,z)=(1,0,-1)$.
    \item $\left(7,\frac{13}{2},\sqrt{10}\right)$, $\left(7,\frac{13}{2},-\sqrt{10}\right)$, $\left(13,\frac72,\sqrt{10}\right)$ arba $\left(13,\frac72,-\sqrt{10}\right)$.

\end{enumerate}


\subsection{Keitiniai}

\begin{enumerate}
    \renewcommand{\labelenumi}{\textbf{\theenumi.}}

    \item $A=t^3$.
    \item $x+y=1$ arba $x+y=-2$.
    \item $x=-9\pm\sqrt{61}$.
    \item $x=\pm\frac12$.
    \item $y\in\{-6,-3,-2,1\}$.
    \item $x=0$.
    \item $x=-4\pm3\sqrt2$.
    \item $x=\frac{1+\sqrt5}{2}$.
    \item $x=1\pm\sqrt7$ arba $x=\frac{13\pm\sqrt{193}}{2}$.
    \item $x=-1$ arba $x=1$.
    \item $x=1$ arba $x=4$.
    \item $x=\frac{1+\sqrt{13}}{4}$ arba $x=9+6\sqrt3$.
    \item $x=5\pm\sqrt{21}$ arba $x=\frac{-9\pm\sqrt{65}}{2}$.
    \item $x=49$.
    \item $x=13\pm\sqrt3$.

\end{enumerate}

\end{document}